\section{Lecture 2, Part 1: Structured Outputs}

\begin{frame}
\vfill
\centering
{\color{RLBlue}\large Lecture 2}

\vspace{0.3em}
{\color{RLBlue}\Large\bfseries Part 1}

\vspace{0.6em}
{\Large Structured Outputs}
\vfill
\end{frame}

\begin{frame}{Learning Goals}
\begin{itemize}
  \item Say why ``reply in JSON'' is a probability rather than a guarantee, and price the difference.
  \pause
  \item Explain how constrained decoding makes an invalid token unsamplable, and what an automaton index buys (Willard \& Louf, 2023).
  \pause
  \item Name precisely what constraining does \emph{not} fix.
\end{itemize}
\end{frame}

% ─── Motivation ─────────────────────────────────────────────────────────────
\begin{frame}{Why Isn't ``Reply in JSON'' Enough?}
\begin{itemize}
  \item Asking politely works most of the time. \emph{Most of the time} is exactly the problem.
  \pause
  \item A trajectory is not one call but $n$. If each output is well-formed with probability $q$, the whole run parses with $q^{n}$ -- the same multiplicative arithmetic that governed reliability in Lecture 1.
  \pause
  \item \textbf{Numerical:} $q = 0.98$ sounds like a solved problem. Over a $10$-step trajectory, $0.98^{10} = 0.82$ -- about one run in five dies on formatting alone, having done nothing wrong.
  \pause
  \item And each failure is not free: a retry re-sends the whole prefix, at the uncached rate, for a turn that produced no information.
\end{itemize}
\end{frame}

\begin{frame}{Three Levels of Enforcement}
\begin{itemize}
  \item \textbf{Ask.} Put the schema in the prompt and hope. Zero engineering; fails silently and at a rate you are not measuring.
  \pause
  \item \textbf{Validate and retry.} Parse the output against the schema; on failure, hand the validation error back as an observation and let the agent correct itself. Recovers -- but pays a full extra turn each time.
  \pause
  \item \textbf{Constrain.} Make an invalid token impossible to sample in the first place. There is no failure to recover from, so there is no retry.
  \pause
  \item These compose rather than compete: \textbf{constrain the syntax, validate the semantics}. No decoder can check that article $182$ exists; only your code can.
\end{itemize}
\end{frame}

\begin{frame}{Constrained Decoding}
\begin{rldefinition}
\textbf{Constrained decoding} -- before sampling each token, set the score of every token that cannot continue a valid output to $-\infty$. \emph{Here:} directly after \texttt{\{"article\_id": "} only digit tokens survive.

\smallskip
\textbf{Mask} -- the vector over the vocabulary saying which tokens are still legal. \emph{Here:} it is recomputed at every step, because what is legal depends on what has already been emitted.

\smallskip
\textbf{Guide} -- the thing the mask is derived from: a regular expression, or a grammar. \emph{Here:} the JSON Schema from Part~0 is compiled into one automatically.
\end{rldefinition}
\pause
\begin{itemize}
  \item Nothing about the model changes. This is sampling-time surgery on the distribution the model already produced.
\end{itemize}
\end{frame}

\begin{frame}{Worked Example: Masking One Token}
\small
The model has emitted
\begin{center}
\texttt{\{"article\_id": "18}
\end{center}
and is choosing the next token. The guide says: inside a digit string, either another digit or the closing quote. Tokens that begin with a space are shown indented, as the tokenizer stores them.
\vspace{0.4em}

\centering
\small
\begin{tabular}{l|c|l}
Candidate token & Model's score & After masking \\
\hline
\texttt{2}        & $-0.4$ & $-0.4$ \\
\texttt{"}        & $-1.9$ & $-1.9$ \\
\texttt{\ the}    & $-3.1$ & $-\infty$ \\
\texttt{\}}       & $-4.0$ & $-\infty$ \\
\texttt{\ Article}& $-4.6$ & $-\infty$ \\
\end{tabular}

\vspace{0.5em}
\pause
\raggedright
\small
Sampling then happens over the survivors only. The model does not become \emph{unlikely} to write prose in the middle of an identifier -- it becomes \textbf{unable} to.
\end{frame}

% ─── Willard & Louf ─────────────────────────────────────────────────────────
\begin{frame}{Regular Expressions, Automata, and an Index}
\begin{columns}[T]
\column{0.48\textwidth}
\small
\begin{itemize}
  \item Willard \& Louf (2023) reframe guided generation as \textbf{walking a finite-state machine}: what has been emitted so far is a state, and each vocabulary token either is a legal edge out of it or is not.
  \pause
  \item \textbf{The naive cost:} test every token in the vocabulary against the guide, at every step. Vocabularies run past $100{,}000$ tokens.
\end{itemize}
\column{0.50\textwidth}
\centering
\begin{tikzpicture}[
  font=\scriptsize,
  st/.style={draw=RLBlue, fill=RLGray, circle, inner sep=1.2pt, minimum size=6.5mm},
  acc/.style={draw=RLGreen, fill=RLGray, circle, double, inner sep=1.2pt,
              minimum size=6.5mm},
  lbl/.style={font=\scriptsize, inner sep=1.5pt}
]
  \node[st]  (a) at (0,0)    {$q_0$};
  \node[st]  (b) at (1.55,0) {$q_1$};
  \node[st]  (c) at (3.1,0)  {$q_2$};
  \node[acc] (d) at (4.65,0) {$q_3$};
  \draw[->] (a) -- node[lbl,above] {\texttt{"}} (b);
  \draw[->] (b) -- node[lbl,above] {digit} (c);
  \draw[->] (c) -- node[lbl,above] {\texttt{"}} (d);
  \draw[->] (c) to[loop below] node[lbl,below] {digit} (c);
\end{tikzpicture}

\vspace{0.4em}
\scriptsize The guide for an article identifier -- at least one digit, then close. In $q_2$ the legal set is \{digits, \texttt{"}\}, and it is the same set every time the walk returns to $q_2$.
\end{columns}
\pause
\vspace{0.3em}
\begin{itemize}
  \small
  \item \textbf{Their construction:} because the legal set depends only on the \emph{state}, pre-compute it once per state and store an index. At generation time the mask is a lookup, so the overhead stops growing with vocabulary size -- guided generation runs at close to unguided speed. Implemented in the \texttt{outlines} library.
\end{itemize}
\end{frame}

\begin{frame}{Grammars, Because JSON Is Not Regular}
\begin{itemize}
  \item A finite-state machine cannot count. JSON nests to arbitrary depth -- \texttt{\{"a": \{"b": [1, 2]\}\}} -- and matching every brace to its partner needs memory that a fixed set of states does not have.
  \pause
  \item So the guide is a \textbf{context-free grammar}, and the machine that walks it is a pushdown automaton: the same state machine, plus a stack recording which brackets are still open.
  \pause
  \item The indexing idea survives intact. The state is now the pair (machine state, stack top), and the allowed-token set is still a lookup from it.
  \pause
  \item This is what \textbf{grammar-guided sampling} means in practice -- \texttt{llama.cpp}'s GBNF and \texttt{outlines} both take a grammar and hand back a mask.
\end{itemize}
\end{frame}

\begin{frame}{What Constraining Does Not Buy You}
\begin{rlpitfall}
Constrained decoding guarantees \textbf{shape}, never \textbf{truth}. \texttt{\{"article\_id": "999", "quote": "..."\}} is flawless JSON, matches the schema perfectly, and cites an article that does not exist.
\end{rlpitfall}
\pause
\begin{itemize}
  \item \textbf{What you can still push into the type:} make \texttt{article\_id} an enum over the identifiers that actually exist. A nonexistent citation then becomes unsamplable too -- this is the enum rule from Part~0, applied where it pays most.
  \pause
  \item \textbf{What no decoder can reach:} whether the quoted sentence really appears in that article, and whether it supports the claim it is attached to. That is verification, not decoding.
  \pause
  \item It arrives twice more in this module: as the verifier subagent in the capstone's Sprint~3, and as the groundedness measure in Lecture 8.
\end{itemize}
\end{frame}

\begin{frame}{Why This Matters: Structured Citations Are Sprint 1}
\begin{itemize}
  \item The capstone's first deliverable is an agent whose every claim carries a machine-checkable citation. That is a schema, not a prompt instruction.
  \pause
  \item \textbf{Constrain the citation object:} \texttt{claim} a string, \texttt{quote} a string, \texttt{article\_id} an enum over the $7{,}368$ article identifiers in the corpus.
  \pause
  \item ``Did it cite something that exists?'' then stops being a measurement and becomes a guarantee -- and your evaluation budget moves to the question that actually matters: \emph{does the quote support the claim?}
  \pause
  \item Magesh et al.\ (2024) found $17$--$33\%$ hallucination in commercial tools that were already doing retrieval. Shape is the cheap half of that problem. Support is the expensive half.
\end{itemize}
\end{frame}

\begin{frame}{Summary}
\begin{itemize}
  \item A \textbf{tool schema} turns a request the model wrote in prose into a name and typed arguments your code can dispatch on -- and it is the model's only documentation.
  \pause
  \item Push every closed set into the type. \textbf{Enums beat sentences}, exactly as arguments beat extra tools in Lecture 1.
  \pause
  \item Models come to call tools by being \textbf{told} (prompt), \textbf{trained} (Toolformer), or \textbf{trained around a retriever} (Gorilla) -- the last because APIs change after training day.
  \pause
  \item \textbf{Constrained decoding} masks illegal tokens before sampling; an index over automaton states makes the mask a lookup rather than a scan (Willard \& Louf, 2023).
  \pause
  \item Constrain the syntax, \textbf{validate the semantics}. Valid JSON is not a true citation.
\end{itemize}
\end{frame}

\begin{frame}{Reading}
\begin{itemize}
  \item Patil, Zhang, Wang \& Gonzalez (2023). \textit{Gorilla: Large Language Model Connected with Massive APIs.} --- APIBench, retriever-aware training, and AST-based hallucination scoring.
  \pause
  \item Willard \& Louf (2023). \textit{Efficient Guided Generation for Large Language Models.} --- guided generation as an automaton walk, and the vocabulary index that makes it cheap.
  \pause
  \item \textit{JSON Schema} and \textit{OpenAPI} specifications --- the two formats every provider's tool-calling interface is built on. Read the \texttt{enum}, \texttt{required} and \texttt{description} keywords properly; they are the whole lecture.
  \pause
  \item \textbf{Next lecture:} the trace stops fitting. Context limits, compression, and what an agent should remember once the scratchpad is full.
\end{itemize}
\end{frame}
